\documentclass[11pt]{article}

\usepackage[margin=1in]{geometry}
\usepackage[utf8]{inputenc}
\usepackage[T1]{fontenc}
\usepackage{amsmath, amssymb}
\usepackage{booktabs}
\usepackage{enumitem}
\usepackage{titlesec}
\usepackage{fancyhdr}
\usepackage{float}
\usepackage{xcolor}
\usepackage{array}
\usepackage{listings}

\usepackage[hidelinks]{hyperref}

\titlespacing*{\section}{0pt}{1.2em}{0.5em}
\titlespacing*{\subsection}{0pt}{0.8em}{0.4em}
\titleformat{\section}{\large\bfseries}{\thesection.}{0.5em}{}
\titleformat{\subsection}{\normalsize\bfseries}{\thesubsection}{0.5em}{}

\pagestyle{fancy}
\fancyhf{}
\lhead{\small CS 388: Natural Language Processing}
\rhead{\small Homework 4}
\cfoot{\small\thepage}
\renewcommand{\headrulewidth}{0.4pt}

\setlist[enumerate]{itemsep=2pt, topsep=3pt}
\setlist[itemize]{itemsep=2pt, topsep=3pt}

\newcommand{\wrd}[1]{\texttt{#1}}

\definecolor{solblue}{RGB}{0,60,180}
\newcommand{\sol}[1]{\par\medskip{\color{solblue}\noindent\textbf{Solution.} #1}\par\medskip}
\newcommand{\B}[1]{{\color{solblue}#1}}

% --- code listing style ---
\lstdefinestyle{py}{
  language=Python,
  basicstyle=\ttfamily\footnotesize,
  commentstyle=\itshape,
  keywordstyle=\bfseries,
  numbers=left,
  numberstyle=\tiny,
  numbersep=8pt,
  frame=single,
  framesep=5pt,
  xleftmargin=18pt,
  showstringspaces=false,
  columns=fullflexible,
  keepspaces=true,
}

\begin{document}

\begin{center}
  {\Large\bfseries Homework 4: Reasoning, Vision-Language Models, and Agents}\\[2pt]
  {\normalsize CS 388: Natural Language Processing \quad\textbullet\quad Fall 2026}\\[2pt]
\end{center}
\vspace{0.5em}
\hrule
\vspace{1em}

\noindent
\textbf{Instructions.} Please transcribe your answers/work into the provided \LaTeX file and submit
the PDF.
\textbf{Please use the {\tt{\textbackslash B\{\}}} command to wrap your solutions, making them blue (this will help us when grading).}

\section{Question 1: RLVR and GRPO (35 points)}

Below is one training step of a simplified version of GRPO for teaching a model to solve math word problems with
reasoning. For simplicity, it leaves out the ratio clipping and the KL penalty to a reference model; you do not need
to add them. The verifier is supposed to give reward $1$ if the \emph{final answer} of a completion is correct and
$0$ otherwise.

\medskip\noindent
There are \textbf{two} logical mistakes in this code: places where the algorithm deviates from GRPO as we saw it
in class. The code runs without errors. The errors are
logical errors, as opposed to coding errors: you can assume, for example, that the helper functions
({\tt{generate}}, {\tt{extract\_answer}}, {\tt{log\_prob}}) do what their comments say, and that \texttt{mean} and
\texttt{std} compute the mean and standard deviation of a list.

\begin{lstlisting}[style=py]
def grpo_step(policy, optimizer, batch, G=8):
    # batch: list of (question, answer) pairs, e.g. ("5 pens cost $10. What do 3 cost?", "6")
    optimizer.zero_grad()
    total_loss = 0.0

    for question, answer in batch:
        # sample a group of G completions (reasoning followed by a final answer)
        completions = [policy.generate(question, temperature=1.0) for _ in range(G)]
        # extract_answer(c) returns the final answer at the end of completion c

        # verifier
        rewards = [float(answer in c) for c in completions]

        # group-relative advantages
        m, s = mean(rewards), std(rewards)
        advs = [(r - m) / (s + 1e-8) for r in rewards]

        for c, adv in zip(completions, advs):
            # log_prob(y, given=x): sum of the log-probs of the tokens of y given x
            logp = policy.log_prob(answer, given=question)
            total_loss = total_loss - adv * logp

    (total_loss / (len(batch) * G)).backward()
    optimizer.step()
\end{lstlisting}

\begin{enumerate}[label=(\alph*)]

  \item \textbf{(5 pts)} Find the first mistake. Give the line number(s) and a fix (pseudocode is fine).

  \item \textbf{(7 pts)} Explain the effect of this mistake. What would happen if we failed to fix it (ignoring
  the other error)? Describe concretely what kind of completions the model would learn to produce.

  \item \textbf{(5 pts)} Find the second mistake. Give the line number(s) and a fix.

  \item \textbf{(7 pts)} Explain the effect of this mistake. What would happen if we failed to fix it (ignoring
  the other error)? \emph{Hint:} look at what the loss for a single question adds up to.

  \item \textbf{(6 pts)} Now assume both mistakes are fixed. For a question where all $G = 8$ completions are
  wrong, what are the advantages, and how much does this question contribute to the gradient? What about a
  question where all $8$ completions are correct?

  \item \textbf{(5 pts)} You start GRPO training from a base model that solves essentially $0\%$ of the training
  questions (none of its sampled completions are ever correct). Based on (e), what happens during training?
  Suggest one way to fix this.

\end{enumerate}

\section{Question 2: Vision-Language Models (30 points)}

\begin{enumerate}[label=(\alph*)]

  \item \textbf{(12 pts)} \textbf{Designing a VLM application.} A city runs a 311 app where residents can submit a
  photo of a road problem. The city wants to automatically (1) decide whether the photo actually shows a pothole,
  (2) rate its severity as \emph{minor}, \emph{moderate}, or \emph{severe}, and (3) write a one-sentence description
  for the repair crew (e.g.\ \emph{``Large pothole in the right lane, next to a storm drain''}). Describe how you
  would address this with a VLM. Your answer should cover (3 pts each):
  \begin{enumerate}[label=(\roman*)]
    \item what the model's input and output look like, including the prompt and the output format;
    \item whether you would use the VLM zero-shot (prompting only) or fine-tune it, and what data you would need;
    \item how you would evaluate the system before deploying it;
    \item one realistic failure mode, and how you would mitigate it.
  \end{enumerate}

  \item \textbf{(18 pts)} \textbf{Reviewing an architecture.} You are running an engineering team building a VLM
  for pharmacies. A pharmacist places pills on a counting tray and takes a photo from above with a fixed camera;
  the model must answer \emph{``How many pills are on the tray?''} The photos are $2048 \times 2048$ pixels, each
  pill is roughly $40 \times 40$ pixels in the photo, and a tray holds up to $200$ pills, often touching each other.

  \medskip
  Engineer B proposes the following model:
  \begin{center}
  \small
  \begin{tabular}{>{\raggedright\arraybackslash}p{4.4cm} >{\raggedright\arraybackslash}p{9.6cm}}
  \toprule
  Component & Setting\\
  \midrule
  Image preprocessing & Resize each photo to $448 \times 448$ pixels\\
  Vision encoder & ViT with patch size $64 \times 64$; $24$ layers, hidden size $1024$, $16$ attention heads;
  pretrained with a CLIP-style objective on $400$M image--caption pairs\\
  Projector & 2-layer MLP mapping each patch embedding to one LLM input token\\
  Language model & $7$B-parameter decoder-only LLM ($32$ layers, hidden size $4096$), context length $4096$\\
  Fine-tuning data & $50$k tray photos with the correct count, formatted as question--answer pairs\\
  Optimization & AdamW, learning rate $2 \times 10^{-5}$, batch size $256$, $3$ epochs\\
  \bottomrule
  \end{tabular}
  \end{center}

  \begin{enumerate}[label=(\roman*)]
    \item \textbf{(6 pts)} How many visual tokens does the LLM receive for each photo? How many pixels of the
    \emph{original} photo does one patch cover (along each side), and roughly how many pills could fall inside a
    single patch? Show your work.
    \item \textbf{(6 pts)} Do you let Engineer B train this model? Why or why not?
    \item \textbf{(6 pts)} If not, what would you change? Be specific (give numbers), and describe what your change
    costs: how many visual tokens does your design produce, and does it fit in the LLM's context?
  \end{enumerate}

\end{enumerate}

\section{Question 3: Agents (35 points)}

A user asks an LLM agent:
\begin{quote}
\emph{``Schedule a one-hour meeting next week with Alice, Bob, and Carol.''}
\end{quote}
The agent can read the user's calendar and the calendars of Alice and Bob, which are shared with the user. Carol's
calendar is private. The agent can also look up Carol's \emph{free/busy status} for a given time, send emails to
any of the attendees and read their replies, and send a calendar invitation. The episode ends when the agent sends
an invitation or gives up. The agent receives a reward of $+1$ if it sends an invitation for a time when all four
people can attend, $-1$ if it sends an invitation for a time when someone cannot attend (the meeting will have to
be rescheduled), and $0$ if it sends no invitation.

\begin{enumerate}[label=(\alph*)]

  \item \textbf{(10 pts)} Write this task as a POMDP. Describe each component (states $S$, actions $A$,
  observations $O$, transitions $T(s' \mid s, a)$, and reward $R$) in a sentence or two, specific to this task.

  \item \textbf{(6 pts)} Name two parts of the state that the agent cannot observe directly. For each, name an
  action that gives the agent information about it, and say whether that information is complete or only partial.

  \item \textbf{(9 pts)} The agent finds that Tuesday at 2 pm is free for the user, Alice, and Bob. Based on what it
  knows, it estimates that Carol is free at that time with probability $p = 0.7$. The agent can either:
  \begin{itemize}
    \item \textbf{Invite now.} It sends the invitation for Tuesday at 2 pm.
    \item \textbf{Email Carol first.} Emailing costs a reward of $c$. Assume Carol answers truthfully. While the
    agent waits for her reply, there is a $0.2$ probability that the slot is taken on someone else's calendar. If
    Carol is free and the slot is still open, the agent sends the invitation; otherwise it gives up.
  \end{itemize}
  \begin{enumerate}[label=(\roman*)]
    \item With $c = 0.1$, compute the expected reward of each option. Which should the agent choose?
    \item What is the largest cost $c$ for which emailing Carol first is still the better choice? (This is the
    \emph{value of information} of Carol's reply.)
  \end{enumerate}

  \item \textbf{(10 pts)} Instead of emailing, the agent first looks up Carol's free/busy status for Tuesday at
  2 pm, which is free of cost. The lookup is imperfect: it only knows about events that Carol has put in her
  calendar. If Carol is actually free, the lookup always shows \emph{free}. If she is actually busy, it shows
  \emph{busy} with probability $0.6$ and \emph{free} with probability $0.4$ (she has not entered the event).
  \begin{enumerate}[label=(\roman*)]
    \item The lookup shows \emph{free}. Starting from the prior $p = 0.7$, compute the agent's updated belief that
    Carol is free.
    \item With this updated belief (and still $c = 0.1$), should the agent invite now or email Carol first? Show the
    expected rewards.
    \item What should the agent believe, and do, if the lookup shows \emph{busy} instead?
  \end{enumerate}

\end{enumerate}

\end{document}
